\documentclass{article}
\usepackage[utf8]{inputenc}
\usepackage{amsmath}
\usepackage{geometry}
\geometry{margin=2cm}
\begin{document}

\section*{Questão 154 --- ENEM 2024 --- Matemática}

Contratos de diversos serviços disponíveis na internet apresentam uma quantidade excessiva de informações, tornando o tempo de leitura desses contratos longo. A tabela apresenta o tempo necessário para a leitura completa do contrato de alguns serviços digitais:

\begin{table}[h!]
\centering
\begin{tabular}{|c|c|}
\hline
\textbf{Tipo de serviço} & \textbf{Tempo necessário para leitura completa do contrato (em minutos)} \\
\hline
A & 36 \\
B & 17 \\
C & 27 \\
D & 13 \\
E & 13 \\
F & 13 \\
\hline
\end{tabular}
\caption{Tempo (min) para leitura dos contratos dos serviços A a F}
\end{table}

\section*{Solução Técnica}

O tempo médio necessário para a leitura completa do contrato dos serviços A a F é obtido pela média aritmética dos tempos indicados:
\[
\text{Média} = \frac{36 + 17 + 27 + 13 + 13 + 13}{6} = \frac{119}{6} \approx 19,8333...
\]
Arredondando para uma casa decimal:
\[
19,8 \text{ minutos}
\]

\section*{Fórmulas Utilizadas}
\[
\text{Média aritmética} = \frac{\sum_{i=1}^n x_i}{n}
\]
onde $x_i$ são os valores dos tempos e $n$ é o número de serviços.

\section*{Evidência Visual Utilizada}

Tabela que fornece o tempo necessário para a leitura completa do contrato para os 6 serviços A a F.

\section*{Explicação Didática}

Para encontrar o tempo médio:
\begin{enumerate}
  \item Identificamos os valores na tabela (36, 17, 27, 13, 13, 13). 
  \item Somamos todos: $36+17+27+13+13+13 = 119$ minutos.
  \item Contamos a quantidade de serviços: 6.
  \item Calculamos a média dividindo a soma pela quantidade: $119/6 = 19,8333...$ minutos.
  \item Finalmente, arredondamos para uma casa decimal: 19,8 minutos.
\end{enumerate}

\textbf{Erros comuns a evitar:}
\begin{itemize}
  \item Confundir média com mediana ou moda.
  \item Esquecer valores na soma ou dividir pela quantidade errada.
  \item Arredondar incorretamente o resultado.
  \item Considerar apenas o menor ou maior valor como a média.
\end{itemize}

\textbf{Dicas para o ensino:}
\begin{itemize}
  \item Explicar o conceito e fórmula da média aritmética.
  \item Mostrar como ler corretamente tabelas.
  \item Explicar a importância do arredondamento.
  \item Diferenciar média, mediana e moda para clareza conceitual.
\end{itemize}

\section*{Distratores e Seus Erros}
\begin{description}
\item[Alternativa A:] Escolhe o menor valor $\boxed{13}$, confundindo valor individual com média.
\item[Alternativa B:] Apresenta valor próximo, mas incorreto por erro de cálculo.
\item[Alternativa D:] Utiliza o valor central ou modal $13$, confundindo com mediana.
\item[Alternativa E:] Confunde soma total ou faz divisão errada, chegando a $23,3$.
\end{description}

\section*{Perfil da Questão}
\begin{itemize}
  \item \textbf{Habilidade principal:} Cálculo de média aritmética.
  \item \textbf{Habilidades secundárias:} Interpretação de tabela e operações básicas.
  \item \textbf{Demanda cognitiva:} Baixa a média.
  \item \textbf{Dificuldade estimada:} Média.
  \item \textbf{Observações:} Requer leitura e compreensão de tabela, operações simples e arredondamento.
\end{itemize}


\end{document}